%%%%%%%%%%%%%%%%%%%%%%%%%%%%%%%%%%%%%%%%%%%%%%%%%%%%%%%%%%%%%%%%%%%%%%
%%%%%%%%%%%%%%%%%%%%%%%%%%%%%%%%%%%%%%%%%%%%%%%%%%%%%%%%%%%%%%%%%%%%%%
%%                                                                  %%
%% EJERCICIO STANDALONE - FORMATO TIPO CÉSAR LLORENTE              %%
%%                                                                  %%
%% Anchura: 9 cm                                                     %%
%% Altura: automática + 0,5 cm (2,5 mm arriba y abajo)              %%
%%                                                                  %%
%%%%%%%%%%%%%%%%%%%%%%%%%%%%%%%%%%%%%%%%%%%%%%%%%%%%%%%%%%%%%%%%%%%%%%
%%%%%%%%%%%%%%%%%%%%%%%%%%%%%%%%%%%%%%%%%%%%%%%%%%%%%%%%%%%%%%%%%%%%%%

%%%%%%%%%%%%%%%%%%%%%%%%%%%%%%%%%%%%%%%%%%%%%%%%%%%%%%%%%%%%%%%%%%%%%
% ESTILO Y MEDIDAS                                                  %
%%%%%%%%%%%%%%%%%%%%%%%%%%%%%%%%%%%%%%%%%%%%%%%%%%%%%%%%%%%%%%%%%%%%%

\documentclass[border=0pt]{standalone}

\setlength{\parindent}{0cm}

%%%%%%%%%%%%%%%%%%%%%%%%%%%%%%%%%%%%%%%%%%%%%%%%%%%%%%%%%%%%%%%%%%%%%
% PAQUETES                                                          %
%%%%%%%%%%%%%%%%%%%%%%%%%%%%%%%%%%%%%%%%%%%%%%%%%%%%%%%%%%%%%%%%%%%%%

\usepackage[utf8x]{inputenc}
\usepackage[T1]{fontenc}
\usepackage[spanish,es-noquoting]{babel}
\usepackage[fleqn]{amsmath}
\usepackage{pxfonts}
\usepackage{graphicx}
\usepackage{lipsum}
\usepackage{calc,fp}
\usepackage[dvipsnames]{xcolor}
\usepackage{pgf}
\usepackage{cancel}
\usepackage{tikz}
\usepackage{array,colortbl,multirow,multicol,booktabs,ctable,eurosym}
\usepackage{pgfplots}
\usetikzlibrary{3d,babel,positioning,arrows,shapes,calc}
\usetikzlibrary{arrows.meta,shapes.geometric}
\usepackage{apacite}
\usepackage{fontawesome5}
\usepackage{mathpazo}

\pgfplotsset{compat=1.18}
\usepgfplotslibrary{fillbetween}

\setlength{\mathindent}{0cm}

%%%%%%%%%%%%%%%%%%%%%%%%%%%%%%%%%%%%%%%%%%%%%%%%%%%%%%%%%%%%%%%%%%%%%
% COLORES                                                           %
%%%%%%%%%%%%%%%%%%%%%%%%%%%%%%%%%%%%%%%%%%%%%%%%%%%%%%%%%%%%%%%%%%%%%

\definecolor{teoria}{cmyk}{0.98,0.5,0,0.35}
\definecolor{tecuacion}{cmyk}{0.98,0.5,0,0.5}
\definecolor{tadorno}{cmyk}{0.60,0.17,0,0.05}
\definecolor{tfondo}{cmyk}{0.30,0.08,0,0.02}
\definecolor{ejercicio}{gray}{0.2}
\definecolor{eadorno}{gray}{0.8}
\definecolor{efondo}{cmyk}{0.12,0,0.18,0}
\definecolor{brickred}{HTML}{B6321C}
\definecolor{info}{gray}{0.7}

\def\te#1{\textcolor{teoria!20}{#1}}
\def\ta#1{\textcolor{tadorno}{#1}}
\def\tf#1{\textcolor{tfondo}{#1}}
\def\ex#1{\textcolor{ejercicio}{#1}}
\def\ea#1{\textcolor{eadorno}{#1}}
\def\ef#1{\colorbox{efondo}{#1}}

%%%%%%%%%%%%%%%%%%%%%%%%%%%%%%%%%%%%%%%%%%%%%%%%%%%%%%%%%%%%%%%%%%%%%
% COMANDOS                                                          %
%%%%%%%%%%%%%%%%%%%%%%%%%%%%%%%%%%%%%%%%%%%%%%%%%%%%%%%%%%%%%%%%%%%%%

\DeclareMathOperator{\Det}{det}
\DeclareMathOperator{\rango}{rango}
\DeclareMathOperator{\rg}{rg}
\DeclareMathOperator{\Proy}{proy}
\DeclareMathOperator{\cotan}{cotan}
\DeclareMathOperator{\arccotan}{arc\,cotan}
\DeclareMathOperator{\dist}{dist}
\DeclareMathOperator{\dom}{dom}

\def\d#1{\dom(#1)}
\def\D#1{\dom\left(#1\right)}
\def\R{\mathbb{R}}

\newdimen\linea

\newenvironment{definicion}[1]%
	{\par
	 \vspace{2mm}
	 \large
	 \ta{\rule{3mm}{3mm}}
	 \textcolor{teoria}{\textbf{\textsf{#1}}}
	 \settowidth{\linea}{\ \rule{3mm}{3mm}\textbf{\textsf{#1}}\ }
	 \ta{\rule{\columnwidth-\linea}{3mm}}\par
	 \normalsize
	 \vspace{1mm}}
	{\hfill
	 \ta{\rule{2mm}{2mm}}\ \ta{\rule{2mm}{2mm}}
	 \vspace{1mm}
	 \par}

\def\n#1{\textcolor{black}{#1}}
\def\t#1{\textcolor{teoria}{#1}}

\def\TT#1{
	\textcolor{tadorno}{\rule{9cm}{0.3mm}}
	\par
	\vspace{-1mm}
	\textcolor{tadorno}{\rule{1mm}{3.5mm}}
	\textcolor{teoria}{\textbf{\noindent #1}}
}

\def\f#1{
	\par
	\vspace{1mm}
	\hspace{1cm}
	\textcolor{tecuacion}{#1}
	\par
	\vspace{1mm}
}

\def\ft#1{\hfil\textcolor{tecuacion}{#1}\hfil}

\def\vs{\vspace{2mm}}

\def\ecuacion#1#2{
	\par
	\vspace{2mm}
	\fcolorbox{tadorno}{white}{
		\hbox to 8.65cm{
			\hspace{0.4cm}
			\textcolor{tecuacion}{#1\hfill \textit{#2}}
		}
	}
	\vspace{2mm}
	\par
}

\def\ecuaciond#1#2{
	\par
	\vspace{2mm}
	\fcolorbox{tadorno}{white}{
		\hbox{
			\vbox{
				\hbox{
					\hspace{0.7cm}
					\textcolor{tecuacion}{#1}
				}
				\par
				\hbox to 8.1cm{
					\hfill
					\textcolor{tecuacion}{\textit{#2}}
				}
			}
		}
	}
	\vspace{2mm}
	\par
}

%%%%%%%%%%%%%%%%%%%%%%%%%%%%%%%%%%%%%%%%%%%%%%%%%%%%%%%%%%%%%%%%%%%%%
% COMANDOS GRÁFICOS                                                 %
%%%%%%%%%%%%%%%%%%%%%%%%%%%%%%%%%%%%%%%%%%%%%%%%%%%%%%%%%%%%%%%%%%%%%

% ESTILOS %%%%%%%%%%%%%%%%%%%%%%%%%%%%%%%%%%%%%%%%%%%%%%%%%%%%%%%%%%%

\tikzset{
	trazo teoria/.style={
		color=teoria!50,
		thick
	}
}

\tikzset{
	cuadricula fina/.style={
		color=black!30,
		very thin
	}
}

\tikzset{
	cuadricula normal/.style={
		color=black!30
	}
}

\tikzset{
	cuadricula gruesa/.style={
		color=black!30,
		thick
	}
}

\tikzset{
	normal/.style={
		rectangle,
		inner sep=0pt,
		align=justify
	}
}

\tikzset{
	cuadro/.style={
		rectangle,
		draw=red,
		inner sep=0pt,
		align=justify
	}
}

\def\cuadricula#1{
	\draw[step=0.25cm,cuadricula fina]
		(0,0) grid (9cm,#1);

	\draw[step=0.5cm,cuadricula normal]
		(0,0) grid (9cm,#1);

	\draw[step=1cm,cuadricula gruesa]
		(0,0) grid (9cm,#1);
}

\def\aclarat#1{
	\textcolor{teoria}{\small #1}
}

\def\nt#1#2{
	\tikz[remember picture]{
		\node[
			rounded corners=3pt,
			draw=teoria!50,
			fill=teoria!20,
			thick
		]
		(#1)
		{\textcolor{black}{\small #2}};
	}
}

\newenvironment{Tabular}[2][1]
	{\def\arraystretch{#1}\tabular{#2}}
	{\endtabular}

\def\T#1{
	\hspace{-1.3mm}
	\raisebox{-5mm}[0pt][0pt]{
		\tikz{
			\draw[very thick,tadorno]
				(0.1,0.4) --
				(0.1,0) --
				(0,0) --
				(0,0.4) --
				(9cm,0.4);

			\fill[tadorno]
				(0,0) rectangle (0.1,0.4);
		}
	}
	\\
	\phantom{I}
	\textcolor{teoria}{\textbf{#1}}
}

\def\seccion#1{
	\par
	\vspace{5mm}

	\textcolor{tfondo}{
		\rule{9cm}{8mm}
	}

	\par
	\vspace{-4mm}

	\textbf{
		\Huge
		\phantom{I}
		\textcolor{teoria}{#1}
	}

	\par
	\vspace{3mm}
}

%%%%%%%%%%%%%%%%%%%%%%%%%%%%%%%%%%%%%%%%%%%%%%%%%%%%%%%%%%%%%%%%%%%%%
% CONTADORES                                                        %
%%%%%%%%%%%%%%%%%%%%%%%%%%%%%%%%%%%%%%%%%%%%%%%%%%%%%%%%%%%%%%%%%%%%%

\newcounter{punto}
\newcounter{ejercicio}[punto]
\newcounter{apartado}[ejercicio]

\def\theejercicio{\arabic{ejercicio}.}
\def\theapartado{\theejercicio\alph{apartado}.}

%%%%%%%%%%%%%%%%%%%%%%%%%%%%%%%%%%%%%%%%%%%%%%%%%%%%%%%%%%%%%%%%%%%%%
% ENTORNO STANDALONE                                                %
%%%%%%%%%%%%%%%%%%%%%%%%%%%%%%%%%%%%%%%%%%%%%%%%%%%%%%%%%%%%%%%%%%%%%

% El PDF tendrá exactamente 9 cm de anchura.
%
% La altura se adapta automáticamente al contenido.
%
% Se añaden:
%
%     2,5 mm arriba
%     2,5 mm abajo
%
% Es decir, 0,5 cm adicionales de altura total.

\newenvironment{ejercicios}
{
	\noindent
	\hspace*{2.5mm}
	\begin{minipage}{8.5cm}

	\vspace*{2.5mm}

	\setcounter{ejercicio}{0}
	\setcounter{apartado}{0}
	\renewcommand{\ej}{\refstepcounter{ejercicio}}

	% Conserva el mismo nivel de anidamiento que la plantilla original.
	% La lista es invisible y no consume anchura, pero hace que los
	% apartados usen los parámetros de una lista de segundo nivel.
	\begin{list}{}{
		\setlength{\leftmargin}{0pt}
		\setlength{\rightmargin}{0pt}
		\setlength{\labelwidth}{0pt}
		\setlength{\labelsep}{0pt}
		\setlength{\itemindent}{0pt}
		\setlength{\listparindent}{0pt}
		\setlength{\topsep}{0pt}
		\setlength{\partopsep}{0pt}
		\setlength{\parsep}{0pt}
		\setlength{\itemsep}{0pt}
	}
	\item[]
}
{
	\end{list}

	\par
	\vspace*{2.5mm}

	\end{minipage}
	\hspace*{2.5mm}
}

%%%%%%%%%%%%%%%%%%%%%%%%%%%%%%%%%%%%%%%%%%%%%%%%%%%%%%%%%%%%%%%%%%%%%
% APARTADOS                                                         %
%%%%%%%%%%%%%%%%%%%%%%%%%%%%%%%%%%%%%%%%%%%%%%%%%%%%%%%%%%%%%%%%%%%%%

\newenvironment{apartados}
{
	\begin{list}{\ex{\textbf{\textsf{\alph{apartado}.\hspace{2mm}}}}}{}
}
{
	\end{list}
}

\newenvironment{apartadosc}
{
	\begin{multicols}{2}

	\begin{list}{\ex{\textbf{\textsf{\alph{apartado}.\hspace{2mm}}}}}{}
}
{
	\end{list}

	\end{multicols}
}

\newenvironment{apartadoscc}
{
	\begin{multicols}{3}

	\begin{list}{\ex{\textbf{\textsf{\alph{apartado}.\hspace{2mm}}}}}{}
}
{
	\end{list}

	\end{multicols}
}

\newenvironment{apartadosccc}
{
	\begin{multicols}{4}

	\begin{list}{\ex{\textbf{\textsf{\alph{apartado}.\hspace{2mm}}}}}{}
}
{
	\end{list}

	\end{multicols}
}

%%%%%%%%%%%%%%%%%%%%%%%%%%%%%%%%%%%%%%%%%%%%%%%%%%%%%%%%%%%%%%%%%%%%%
% COMANDOS DE EJERCICIOS                                            %
%%%%%%%%%%%%%%%%%%%%%%%%%%%%%%%%%%%%%%%%%%%%%%%%%%%%%%%%%%%%%%%%%%%%%

\def\ej{\refstepcounter{ejercicio}\item}

\def\ejs{
	$\stackrel{H}{\text{\refstepcounter{ejercicio}}}$
	\item
}

\def\ejss{
	\refstepcounter{ejercicio}
	\item
	\hspace{-14.6mm}
	$\bigstar\bigstar$
	\hspace{7.4mm}
}

\def\ap{\refstepcounter{apartado}\item}

\def\apv{
	\vspace{1mm}
	\refstepcounter{apartado}
	\item
}

%%%%%%%%%%%%%%%%%%%%%%%%%%%%%%%%%%%%%%%%%%%%%%%%%%%%%%%%%%%%%%%%%%%%%
% TÍTULOS                                                           %
%%%%%%%%%%%%%%%%%%%%%%%%%%%%%%%%%%%%%%%%%%%%%%%%%%%%%%%%%%%%%%%%%%%%%

\def\titulos#1{
	\hspace{-2.3mm}

	\raisebox{-5mm}[0pt][0pt]{
		\tikz{
			\draw[very thick,eadorno]
				(0.1,0.4) --
				(0.1,0) --
				(0,0) --
				(0,0.4) --
				(9cm,0.4);

			\fill[eadorno]
				(0,0) rectangle (0.1,0.4);
		}
	}

	\\
	\phantom{I}
	\textcolor{ejercicio}{
		\textbf{#1}
	}
	\par
}

\def\tipo#1{
	\item[]

	\hspace{-8mm}
	\vspace{2mm}

	\tikz{
		\draw[very thick,eadorno]
			(0.1,0.4) --
			(0.1,0) --
			(0,0) --
			(0,0.4) --
			(9cm,0.4);

		\fill[eadorno]
			(0,0) rectangle (0.1,0.4);

		\node at (0.1,0.13)
			[right]
			{\textsf{\small #1}};
	}

	\par
	\vspace{-2mm}
}

%%%%%%%%%%%%%%%%%%%%%%%%%%%%%%%%%%%%%%%%%%%%%%%%%%%%%%%%%%%%%%%%%%%%%
% SOLUCIONES                                                        %
%%%%%%%%%%%%%%%%%%%%%%%%%%%%%%%%%%%%%%%%%%%%%%%%%%%%%%%%%%%%%%%%%%%%%

\def\thesol{
	\ifnum\value{apartado}=0

		\noexpand\medskip

		\noexpand\textit{
			\noexpand\textsf{
				\arabic{ejercicio}.
			}
		}

	\else

		\ifnum\arabic{apartado}=1
			\noexpand\medskip
		\else
		\fi

		\noexpand\textit{
			\noexpand\textsf{
				\arabic{ejercicio}
				\alph{apartado}.
			}
		}

	\fi
}

\def\thelsol{
	\par

	\ifnum\value{apartado}=0

		\noexpand\medskip

		\noexpand\textit{
			\noexpand\textsf{
				\arabic{ejercicio}.
			}
		}

	\else

		\ifnum\arabic{apartado}=1
			\noexpand\medskip
		\else
		\fi

		\noexpand\textit{
			\noexpand\textsf{
				\arabic{ejercicio}
				\alph{apartado}.
			}
		}

	\fi
}

\def\theesol{
	\par

	\ifnum\value{apartado}=0

		\noexpand\medskip

		\noexpand\textit{
			\noexpand\textsf{
				\arabic{ejercicio}.
			}
		}

	\else

		\ifnum\arabic{apartado}=1
			\noexpand\medskip
		\else
		\fi

		\noexpand\textit{
			\noexpand\textsf{
				\arabic{ejercicio}.
			}
		}

	\fi
}

\def\soluciones{
	\par
	\bigskip
	\bigskip
	\noindent

	\tikz{
		\draw[very thick,eadorno]
			(0.1,0.4) --
			(0.1,0) --
			(0,0) --
			(0,0.4) --
			(9cm,0.4);

		\fill[eadorno]
			(0,0) rectangle (0.1,0.4);

		\node[below right]
			at (0.1,0.4)
			{\textbf{Soluciones}};
	}

	\par
	\parindent=3mm
	\hangindent=3mm
	\small
}

%%%%%%%%%%%%%%%%%%%%%%%%%%%%%%%%%%%%%%%%%%%%%%%%%%%%%%%%%%%%%%%%%%%%%
% ICONOS                                                            %
%%%%%%%%%%%%%%%%%%%%%%%%%%%%%%%%%%%%%%%%%%%%%%%%%%%%%%%%%%%%%%%%%%%%%

\def\d{
	\small
	\faPepperHot
	\hspace{1mm}
}

\def\dd{
	\small
	\faPepperHot

	\small
	\faPepperHot

	\hspace{1mm}
}

\def\cg{
	\small
	\faCalculator
	\hspace{1mm}
}

%%%%%%%%%%%%%%%%%%%%%%%%%%%%%%%%%%%%%%%%%%%%%%%%%%%%%%%%%%%%%%%%%%%%%
% INFORMACIÓN Y PUNTUACIONES                                        %
%%%%%%%%%%%%%%%%%%%%%%%%%%%%%%%%%%%%%%%%%%%%%%%%%%%%%%%%%%%%%%%%%%%%%

\def\M#1{\vspace{3mm}{\small\sffamily\textbf{[Puntuación total: #1]}}\hfill}

\def\m#1{
	\hfill
	\textit{[#1]}
}

\def\y{
	\hspace{3mm}
	\text{y}
	\hspace{3mm}
}

%%%%%%%%%%%%%%%%%%%%%%%%%%%%%%%%%%%%%%%%%%%%%%%%%%%%%%%%%%%%%%%%%%%%%
% ACUMULACIÓN DE SOLUCIONES                                         %
%%%%%%%%%%%%%%%%%%%%%%%%%%%%%%%%%%%%%%%%%%%%%%%%%%%%%%%%%%%%%%%%%%%%%

\makeatletter

\def\sol#1{
	\begingroup

	\edef\@num{\thesol}

	\expandafter\g@addto@macro
		\expandafter\soluciones
		\expandafter{
			\@num
			\ 
			\small #1
			\hspace{3mm}
		}

	\endgroup
}

\def\lsol#1{
	\begingroup

	\edef\@num{\thelsol}

	\expandafter\g@addto@macro
		\expandafter\soluciones
		\expandafter{
			\@num
			\ 
			\small #1
			\hspace{3mm}
		}

	\endgroup
}

\def\csol#1{
	\begingroup

	\edef\@num{\thelsol}

	\expandafter\g@addto@macro
		\expandafter\soluciones
		\expandafter{
			\vspace{1mm}
			\\
			\phantom{
				\thelsol
			\hspace{3mm}
			}
			\small #1
			\hspace{3mm}
		}

	\endgroup
}

\makeatother

\makeatletter

\def\esol#1{
	\addtocounter{apartado}{1}

	\begingroup

	\edef\@num{\theesol}

	\expandafter\g@addto@macro
		\expandafter\soluciones
		\expandafter{
			\@num
			\ 
			\small #1
			\hspace{3mm}
		}

	\endgroup
}

\makeatother

\makeatletter

\def\solpg#1{
	\addtocounter{apartado}{1}

	\begingroup

	\edef\@num{\thesol}

	\expandafter\g@addto@macro
		\expandafter\soluciones
		\expandafter{
			\edef\@num{\thesol}

			\tikz{
				\node{
					\expandafter\noexpand\@num
				};
			}
		}

	\endgroup
}

\makeatother

%%%%%%%%%%%%%%%%%%%%%%%%%%%%%%%%%%%%%%%%%%%%%%%%%%%%%%%%%%%%%%%%%%%%%
% MATRICES Y VECTORES                                               %
%%%%%%%%%%%%%%%%%%%%%%%%%%%%%%%%%%%%%%%%%%%%%%%%%%%%%%%%%%%%%%%%%%%%%

\def\talque{\ |\ }

\newenvironment{matrizp}
	{\left(\begin{smallmatrix}}
	{\end{smallmatrix}\right)}

\newenvironment{detp}
	{\left|\begin{smallmatrix}}
	{\end{smallmatrix}\right|}

\newenvironment{sistemap}
	{\left\{\begin{smallmatrix}}
	{\end{smallmatrix}\right.}

\def\Vec#1{
	\overrightarrow{#1}
}

%%%%%%%%%%%%%%%%%%%%%%%%%%%%%%%%%%%%%%%%%%%%%%%%%%%%%%%%%%%%%%%%%%%%%
% COMANDOS TIKZ                                                     %
%%%%%%%%%%%%%%%%%%%%%%%%%%%%%%%%%%%%%%%%%%%%%%%%%%%%%%%%%%%%%%%%%%%%%

\def\punto#1#2#3#4{
	\coordinate (#1) at (#2);

	\draw[fill]
		(#1) circle (0.05);

	\node at (#1)
		[#4]
		{#3};
}

%%%%%%%%%%%%%%%%%%%%%%%%%%%%%%%%%%%%%%%%%%%%%%%%%%%%%%%%%%%%%%%%%%%%%
% PUNTUACIONES                                                      %
%%%%%%%%%%%%%%%%%%%%%%%%%%%%%%%%%%%%%%%%%%%%%%%%%%%%%%%%%%%%%%%%%%%%%

\def\p#1{\if#11\textsf{\small(1 punto)}\else\textsf{\small(#1 puntos)}\fi}

\def\P#1{
	\textbf{
		\p{#1}
	}
}

%%%%%%%%%%%%%%%%%%%%%%%%%%%%%%%%%%%%%%%%%%%%%%%%%%%%%%%%%%%%%%%%%%%%%
% INFORMACIÓN                                                       %
%%%%%%%%%%%%%%%%%%%%%%%%%%%%%%%%%%%%%%%%%%%%%%%%%%%%%%%%%%%%%%%%%%%%%

\def\info#1{
	\par
	\vspace{-2mm}
	\noindent
	\makebox[\linewidth][r]{
		\textcolor{gray}{
			\tiny{
				\textsf{
					\uppercase{#1}
				}
			}
		}
	}
}

%%%%%%%%%%%%%%%%%%%%%%%%%%%%%%%%%%%%%%%%%%%%%%%%%%%%%%%%%%%%%%%%%%%%%
% OTROS COMANDOS                                                    %
%%%%%%%%%%%%%%%%%%%%%%%%%%%%%%%%%%%%%%%%%%%%%%%%%%%%%%%%%%%%%%%%%%%%%

\def\c#1#2{
	\left(
		\begin{smallmatrix}
			#1\\
			#2
		\end{smallmatrix}
	\right)
}

\def\completar#1{
	\vspace{#1}
	\phantom{.}
}

%%%%%%%%%%%%%%%%%%%%%%%%%%%%%%%%%%%%%%%%%%%%%%%%%%%%%%%%%%%%%%%%%%%%%
% MINI GRÁFICAS                                                     %
%%%%%%%%%%%%%%%%%%%%%%%%%%%%%%%%%%%%%%%%%%%%%%%%%%%%%%%%%%%%%%%%%%%%%

\newcommand{\miniplot}[6]{
	\tikz[scale=#1]{
		\begin{axis}[
			domain=#2:#3,
			samples=25,
			axis lines=middle,
			xmin=#2,
			xmax=#3,
			ymin=#4,
			ymax=#5,
		]

			\addplot[
				thick,
				black,
				name path=curve
			]
			{#6};

		\end{axis}
	}
}

%%%%%%%%%%%%%%%%%%%%%%%%%%%%%%%%%%%%%%%%%%%%%%%%%%%%%%%%%%%%%%%%%%%%%
% SOLUCIONES DESARROLLADAS                                          %
%%%%%%%%%%%%%%%%%%%%%%%%%%%%%%%%%%%%%%%%%%%%%%%%%%%%%%%%%%%%%%%%%%%%%

\definecolor{solcolor}{RGB}{113,109,130}

\newenvironment{solucion}
{
	\vspace{2mm}
	\leavevmode
	\color{solcolor}
}
{
	\renewcommand{\arraystretch}{1}
}

%%%%%%%%%%%%%%%%%%%%%%%%%%%%%%%%%%%%%%%%%%%%%%%%%%%%%%%%%%%%%%%%%%%%%
% IGUALDADES COMENTADAS                                             %
%%%%%%%%%%%%%%%%%%%%%%%%%%%%%%%%%%%%%%%%%%%%%%%%%%%%%%%%%%%%%%%%%%%%%

\newcommand{\beq}[1]{
	\ensuremath{
		\underset{
			\sffamily{
				\text{
					\textbf{
						\tiny #1
					}
				}
			}
		}{=}
	}
}

\newcommand{\aeq}[1]{
	\ensuremath{
		\overset{
			\sffamily{
				\text{
					\textbf{
						\tiny #1
					}
				}
			}
		}{=}
	}
}

\newcommand{\lheq}{
	\ensuremath{
		\underset{
			\sffamily{
				\text{
					\textbf{
						\tiny L'H
					}
				}
			}
		}{=}
	}
}

%%%%%%%%%%%%%%%%%%%%%%%%%%%%%%%%%%%%%%%%%%%%%%%%%%%%%%%%%%%%%%%%%%%%%
% FLECHAS DE CAMBIO DE COLUMNA                                     %
%%%%%%%%%%%%%%%%%%%%%%%%%%%%%%%%%%%%%%%%%%%%%%%%%%%%%%%%%%%%%%%%%%%%%

\newcommand{\ccu}[2]{
	\begin{tikzpicture}[
		baseline=(termino.base)
	]

		\node[
			inner sep=0pt
		]
		(termino)
		{\ensuremath{#2}};

		\pgfresetboundingbox

		\useasboundingbox
			(termino.south east)
			rectangle
			(termino.north west);

		\coordinate (A)
			at (
				[shift={(1pt,-1pt)}]
				termino.south east
			);

		\coordinate (B)
			at (
				[shift={(-2pt,2pt)}]
				termino.north west
			);

		\draw[
			-{Stealth},
			opacity=0.8
		]
			(A) -- (B);

		\node[
			anchor=center,
			font=\tiny,
			scale=0.8
		]
			at ($(A)!1.2!(B)$)
			{\ensuremath{#1}};

	\end{tikzpicture}
}

\newcommand{\ccd}[2]{
	\begin{tikzpicture}[
		baseline=(termino.base)
	]

		\node[
			inner sep=0pt
		]
		(termino)
		{\ensuremath{#2}};

		\pgfresetboundingbox

		\useasboundingbox
			(termino.north west)
			rectangle
			(termino.south east);

		\coordinate (A)
			at (
				[shift={(-1pt,1pt)}]
				termino.north west
			);

		\coordinate (B)
			at (
				[shift={(2pt,-2pt)}]
				termino.south east
			);

		\draw[
			-{Stealth},
			opacity=0.8
		]
			(A) -- (B);

		\node[
			anchor=center,
			font=\tiny,
			scale=0.8
		]
			at ($(A)!1.2!(B)$)
			{\ensuremath{#1}};

	\end{tikzpicture}
}

%%%%%%%%%%%%%%%%%%%%%%%%%%%%%%%%%%%%%%%%%%%%%%%%%%%%%%%%%%%%%%%%%%%%%
% COMENTARIOS PEQUEÑOS                                              %
%%%%%%%%%%%%%%%%%%%%%%%%%%%%%%%%%%%%%%%%%%%%%%%%%%%%%%%%%%%%%%%%%%%%%

\newcommand{\fcomment}[1]{
	\makebox[0pt][l]{
		\scalebox{0.6}{
			\ensuremath{#1}
		}
	}
}

\def\si{
	\text{si}
}
